% !TEX root = chapter-standalone.tex

\section{Strictification}
\label{sec:strictification}


\todotext{write this up}

\subsection{The category~\PosL}

We define a category analogous to~\SetL, but its objects are ``tuple posets''.

Given \SY{posets}~$\posA_1,\ldots, \posA_n$ we define the poset
\begin{equation}
    \label{eq:posetoftuples}
    \cObj{\posA_1,\ldots, \posA_n}\definedas \tup{\cObj{\posAsetn{1},\ldots, \posAsetn{n}}, \posleqof{\cObj{\posA_1,\ldots, \posA_n}}},
\end{equation}
where~$\cObj{\posAsetn{1},\ldots, \posAsetn{n}}$ is a set of tuples, and we use the product order:
\begin{equation}
    \prfdoubleperiod{
        \tup{\posela_1,\ldots, \posela_n}
        \posleqof{\cObj{\posA_1,\ldots, \posA_n}}
        \tup{\poselb_1,\ldots, \poselb_n}
    }{
        \posela_i \posleqof{\posA_i} \poselb_i \text{ for all }~i\setin \makeset{1,\ldots, n}
    }
\end{equation}
\begin{definition}[The category~\PosL]
    \SYNDEF{category of tuple-posets and monotone maps}
    \label{def:PosL}
    The category~\PosL is the \SY{subcategory} of \Pos where the objects are tuple posets, \SY{posets} of the form
    \begin{equation}
        \cObj{\posA_1,\ldots, \posA_n} = \tup{\cObj{\posAsetn{1},\ldots, \posAsetn{n}}, \posleqof{\cObj{\posA_1,\ldots, \posA_n}}}.
    \end{equation}
    % \begin{enumerate}
    %     \item \emph{Objects:} \SY{posets} of the form presented in \cref{eq:posetoftuples}.
    %     \item \emph{Morphisms:}
    %           Monotone functions between \SY{posets} of the form presented in \cref{eq:posetoftuples}.
    %     \item \emph{Composition:}
    %           Composition is the usual composition of functions.
    %     \item \emph{Identities:}
    %           Identity functions.
    % \end{enumerate}
\end{definition}

\begin{lemma}\label{lem:PosL-associative-stacking}
    \PosL is  \SY{associative stacking} using the structure induced by tuple concatenation.
\end{lemma}
\begin{proof}
    Analogously to what we did for \SetL, we can define a multiplication operation $\mtimescatob$ in \PosL.
    \todotext{JL: the definition below is incorrect -- the objects in our category are not posets!}

    Given two objects $\posA = \cObj{\posA_1,\ldots, \posA_n}$ and $\posB = \cObj{\posB_1,\ldots, \posB_n}$, we define
    \begin{equation}
        \label{eq:multposstack}
        % \begin{aligned}
        \cObj{\posA_1,\ldots, \posA_n} \mtimescatob \cObj{\posB_1,\ldots, \posB_n} \definedas
        % \tup{\cObj{\posAsetn{1},\ldots, \posAsetn{m}} \cprod \cObj{\posBsetn{1},\ldots, \posBsetn{n}}, \posleqof{\cObj{\posA_1,\ldots, \posA_m}\cprod \cObj{\posB_1,\ldots, \posB_n}}} \\
        \cObj{\posA_1,\ldots, \posA_n, \posB_1,\ldots, \posB_n}.
    \end{equation}
    %    where the order is defined as
    %    \begin{equation}\label{eq:PosL-order}
    %        \prfdoubleperiod{
    %            \pars{\maketupconcat{\posAnel{1}, \posBnel{1}}}
    %            \ %
    %            \posleqof{\posA\mtimescatob \posB}
    %            \ %
    %            \pars{\maketupconcat{\posAnel{2}, \posBnel{2}}}
    %        }{
    %            \pars{\posAnel{1}\posleqof\posA\posAnel{2}}
    %            \ \booland\ %
    %            \pars{\posBnel{1}\posleqof\posB\posBnel{2}}
    %        }
    %    \end{equation}
    For the multiplication on morphisms, we define
    \begin{equation}
        \prfperiod{
            \mora\colon \posA \mto \posB
        }{
            \quad
        }{
            \mora\colon \posC \mto \posD
        }{
            \defmap{
                \pars{
                    \mora
                    \mtimescatmor
                    \morb
                }
            }{
                \posA\mtimescatob\posC
            }{
                \mto
            }{
                \posB\mtimescatob\posD
            }{
                \maketupconcat{\posAel, \posCel}
            }{
                \maketupconcat{\mora(\posAel), \morb(\posCel)}
            }
        }
    \end{equation}
    We need to check that the expression $\maketupconcat{\mora(\posAel), \morb(\posCel)}$ is monotone.
    This can be easily seen because the order on $\posB\mtimescatob\posD$ is akin to a product order.
    The proof for associativity of $\mtimescatmor$ is the same as in the proof of $\SetL$ (\cref{lem:SetL-is-associative-stacking}).
\end{proof}
% ~$\tup{\posAnel{1},\ldots, \posAnel{m},\posBnel{1},\ldots, \posBnel{n}} \posleqof{\cObj{\posA_1,\ldots, \posA_m}\mtimescatob \cObj{\posB_1,\ldots, \posB_n}} \tup{\posAnel{1}',\ldots, \posAnel{m}',\posBnel{1}',\ldots, \posBnel{n}'}$ if and only if~$\posAnel{i}\posleqof{\posA_i}\posAnel{i}'$ and~$\posBnel{j}\posleqof{\posB_j}\posBnel{j}'$ for all~$i\setin \makeset{1,\ldots, m}$,~$j\setin \makeset{1,\ldots, n}$.
% % Note that this construction is analogous to the product poset.
% The resulting \SY{poset} of tuples is written as~$\cObj{\posA_1,\ldots, \posA_m,\posB_1,\ldots, \posB_n}$.
% We now want to show that one can define a structure for which~\PosL is  \SY{associative stacking}.
% To do so, we define the stacking operation on objects as the multiplication operation~$\cprod$ defined in \cref{eq:multposstack}.

% The stacking operation on morphisms is defined as
% \begin{equation}
%     \defmapcomma{
%         \mtimescatmor
%     }{
%         \Mor_{\PosL} \cartprod \Mor_{\PosL}
%     }{
%         \mto
%     }{
%         \Mor_{\PosL}
%     }{
%         \tup{\mora,\morb}
%     }{
%         \mora \cprod \morb
%     }
% \end{equation}
% where~$(\mora \mtimescatmor \morb)(\tup{\posAnel{1},\ldots, \posAnel{m}} \tupconcat \tup{\posBnel{1},\ldots, \posBnel{n}})\definedas \mora(\tup{\posAnel{1},\ldots, \posAnel{m}})\tupconcat \morb(\tup{\posBnel{1},\ldots, \posBnel{n}})$.

% \begin{lemma}[\PosL is associative stacking]\label{lem:PosL-associative-stacking}
%     \PosL, equipped with the aforementioned stacking operations on objects and morphisms, is  \SY{associative stacking}.
% \end{lemma}

% \begin{proof}
%     % We show that~\PosL is an  \SY{associative stacking} semicategory.
%     We first show that it is a stacking category, by showing that the stacking operations fulfill \cref{def:simple-stacking-semi-cat}, and then showing the compatibility condition.
%     The stacking operation on objects fulfills \cref{def:simple-stacking-semi-cat}, since, given two \SY{posets} of tuples it returns a \SY{poset} of tuples.
%     To show that the stacking operation on morphism fulfills \cref{def:simple-stacking-semi-cat}, we need to show that given monotone~$\mora,\morb$, the map~$\mapa \mtimescatmor \mapb$ is monotone.
%     This is easy to show, by leveraging the monotonicity of~$\mora$ and~$\morb$.
%     For brevity, in the following we write~$\posAnel{}=\tup{\posAnel{1},\ldots, \posAnel{m}}$,~$\posAnel{}'=\tup{\posAnel{1}',\ldots, \posAnel{m}'}$,~$\posBnel{}=\tup{\posBnel{1},\ldots, \posBnel{n}}$, and~$\posBnel{}'=\tup{\posBnel{1}',\ldots, \posBnel{n}'}$.
%     We have
%     \begin{equation}
%         \prfcomma{\posAnel{}\tupconcat \posBnel{} \posleqof{\cObj{\posA_1,\ldots, \posA_m,\posB_1,\ldots, \posB_n}} \posAnel{}'\tupconcat \posBnel{}'}
%         {\prftree{
%                 \left(\posAnel{} \posleqof{\cObj{\posA_1,\ldots, \posA_m}} \posAnel{}'\right) \booland \left(\posBnel{} \posleqof{\cObj{\posB_1,\ldots, \posB_n}} \posBnel{}'\right)}
%             {
%                 \prftree{
%                     \left(\mora(\posAnel{}) \posleqof{\cObj{\posC_1,\ldots, \posC_o}} \mora(\posAnel{}')\right) \booland \left(\morb(\posBnel{})\posleqof{\cObj{\posD_1,\ldots, \posD_t}} \morb(\posBnel{}')\right)}
%                 {
%                     \prftree{
%                         {\mora(\posAnel{}) \tupconcat \morb(\posBnel{})  \posleqof{\cObj{\posC_1,\ldots, \posC_o, \posD_1,\ldots, \posD_t}} \mora(\posAnel{}') \tupconcat \morb(\posBnel{}')}}
%                     {(\mora \mtimescatmor \morb)(\posAnel{}\tupconcat \posBnel{}) \posleqof{\cObj{\posC_1,\ldots, \posC_o,\posD_1,\ldots, \posD_t}} (\mora \mtimescatmor \morb)(\posAnel{}'\tupconcat \posBnel{}')
%                     }}}
%         }
%     \end{equation}
%     which shows monotonicity.
%     Furthermore, the two stacking operations are compatible:
%     \begin{equation}
%         \prfperiod{\mora\colon \cObj{\posA_1,\ldots,\posA_m}\mto \cObj{\posC_1,\ldots,\posC_o}}
%         {\morb\colon \cObj{\posB_1,\ldots,\posB_n}\mto \cObj{\posD_1,\ldots,\posD_t}}
%         {\mora \mtimescatmor \morb \colon \cObj{\posA_1, \ldots, \posA_m, \posB_1, \ldots, \posB_n} \mto \cObj{\posC_1, \ldots, \posC_o, \posD_1, \ldots, \posD_t}}
%     \end{equation}
%     %
%     The proof of associativity is parallel to the proof for~\SetL.
% \end{proof}
%

\subsection{The category~\RelL}
We define a category~\RelL where the objects are sets of tuples (as in~\SetL).

\begin{definition}[The category~\RelL]
    \label{def:RelL}
    \SYNDEF{category of tuple-sets and relations}
    The category~\RelL is the \SY{subcategory} of~\Rel where the objects are tuples sets (objects of \SetL).
    % \begin{enumerate}
    %     \item \emph{Objects:} sets of tuples (as in~\SetL)
    %     \item \emph{Morphisms:}
    %           relations between sets of tuples.
    %     \item \emph{Composition:}
    %           Composition is the usual composition of relations.
    %     \item \emph{Identities:}
    %           The identity morphism on an object~$\cObj{\setAn{1},\ldots, \setAn{m}}$ is given by the identity relation~$\catid_{\cObj{\setAn{1},\ldots, \setAn{m}}}$.
    % \end{enumerate}
\end{definition}

\begin{lemma}\label{lem:RelL-associative-stacking}
    \RelL is  \SY{associative stacking} with the structure induced by tuple concatenation.
\end{lemma}
\begin{proof}
    The multiplication $\mtimescatob$ is $\cprod$, the same as the one defined for $\SetL$.

    The multiplication $\mtimescatmor$ is defined as follows:
    \begin{equation}
        \prf{
            \relA\setsubseteq \setA \cartprod \setB
        }{
            \quad
        }{
            \relB\setsubseteq \setC\cartprod \setD
        }{
            \pars{\relA\mtimescatmor\relB} \setsubseteq
            \pars{\makecprod{\setA,\setC}} \cartprod
            \pars{\makecprod{\setB,\setD}}
        }
    \end{equation}
    where
    \begin{equation}
        \pars{
            \relA
            \mtimescatmor
            \relB
        }
        =
        \makesett{
            \tup{\setAeln{}\tupconcat \setCeln{},\setBeln{}\tupconcat \setDeln{}}
            \mid
            % \setAeln{} \setin \setA,
            % \setBeln{} \setin \setB,
            % \setCeln{} \setin \setC,
            % \setDeln{} \setin \setD:
            (\inrel{\setAeln{}}\relA{\setBeln{}})
            \booland
            (\inrel{\setCeln{}}\relB{\setDeln{}})
        }.
    \end{equation}
    The rest of the proof is left as an exercise.
\end{proof}

% What is a relation in this context?
% Given two sets of tuples~$\cObj{\setAn{1},\ldots, \setAn{m}}$ and~$\cObj{\setBn{1},\ldots, \setBn{n}}$, a relation between them is given by
% \begin{equation}
%     \relA \setsubseteq \cObj{\setAn{1},\ldots, \setAn{m}} \cartprod \cObj{\setBn{1},\ldots, \setBn{n}}.
% % \end{equation}
% We can define a multiplication operation for relations.
% Consider
% \begin{equation}
%     \begin{aligned}
%         \relA\colon \cObj{\setAn{1},\ldots, \setAn{m}} & \mtoin{\RelL}\cObj{\setBn{1},\ldots, \setBn{n}}, \\
%         \relB\colon \cObj{\setCn{1},\ldots, \setCn{o}} & \mtoin{\RelL}\cObj{\setDn{1},\ldots, \setDn{p}}.
%     \end{aligned}
% \end{equation}
% We have
% \begin{equation}
%     \relA\mtimescatmor \relB \colon \cObj{\setAn{1},\ldots, \setAn{m}}\cprod \cObj{\setCn{1},\ldots, \setCn{o}}\mtoin{\RelL}
%     \cObj{\setBn{1},\ldots, \setBn{n}}\cprod \cObj{\setDn{1},\ldots, \setDn{p}}.
% \end{equation}
% Specifically, the multiplication is given by:
% \begin{equation}
%     \label{eq:multiofrels}
%     \defmapperiod{
%         \mtimescatmor
%     }{
%         \Mor_{\RelL}\cartprod \Mor_{\RelL}
%     }{
%         \mto
%     }{
%         \Mor_{\RelL}
%     }{
%         \tup{\relA,\relB}
%     }{
%         \makeset{
%             \tup{\setAeln{}\tupconcat \setCeln{},\setBeln{}\tupconcat \setDeln{}}
%             \mid
%             (\inrel{\setAeln{}}\relA{\setBeln{}}) \booland (\inrel{\setCeln{}}\relB{\setDeln{}})
%         }
%     }
% \end{equation}

\begin{exercise}
    \label{ex:relassocstack}
    Consider the stacking operations on objects and on morphisms introduced in this section.
    Prove that~\RelL is \SY{associative stacking}.
\end{exercise}

\begin{solution}
    To prove the statement we first check that the stacking operations satisfy \cref{def:simple-stacking-semi-cat}, we then show that they are compatible, and finally show associativity.
    The stacking operation on objects was already checked for~\SetL.
    The stacking operation on morphisms clearly returns a valid relation.
    Furthermore, compatibility is satisfied:
    \begin{equation}
        \prfperiod{
            \relA\colon \cObj{\setAn{1},\ldots, \setAn{m}} \mtoin{\RelL}\cObj{\setBn{1},\ldots, \setBn{n}}
        }
        {
            \relB\colon \cObj{\setCn{1},\ldots, \setCn{o}} \mtoin{\RelL}\cObj{\setDn{1},\ldots, \setDn{p}}
        }
        {
            \relA\mtimescatmor \relB \colon \cObj{\setAn{1},\ldots, \setAn{m}}\cprod \cObj{\setCn{1},\ldots, \setCn{o}}\mtoin{\RelL}
            \cObj{\setBn{1},\ldots, \setBn{n}}\cprod \cObj{\setDn{1},\ldots, \setDn{p}}
        }
    \end{equation}
    Finally, associativity for the operation on objects was already shown for~\SetL.
    % The stacking operation on morphisms follows directly from \cref{eq:multiofrels}.
\end{solution}

subsection{The category~$\DPL$}

Analogously, we define the category $\DPL$.

\begin{definition}[The category \DPL]\label{def:DPL}
    \SYNDEF{category of tuple-posets and design problems}
    The category~$\DPL$ is the \SY{subcategory} of \DP where the objects are \SY{posets} of tuple posets (objects of $\PosL$).
    % \begin{enumerate}
    %     \item \emph{Objects:}
    %           \SY{posets} of tuples.
    %     \item \emph{Morphisms:}
    %           A morphism from~$\posA =\cObj{\posA_1,\ldots, \posA_m}$ to ~$\posB = \cObj{\posB_1,\ldots, \posB_m}$ is design problem:
    %           \begin{equation}
    %               \adp\colon \posA\op\cartprod \posB \toinPos \Bool.
    %           \end{equation}
    %     \item \emph{Composition:}
    %           Composition is the usual composition of \SY{design problems}.
    %     \item \emph{Identities:}
    %           The identity morphism on an object~$\posA=\cObj{\posA_1,\ldots,\posA_n}$ is given by the identity \SY{design problem}~$\catid_\posA$.
    % \end{enumerate}
\end{definition}
% \begin{remark}
%     In other words, \DPL is the subcategory of \DP in which objects are of the form presented in \cref{eq:posetoftuples}.
% \end{remark}
\begin{lemma}
    \DPL is \SY{associative stacking} using the structure induced by tuple concatenation.
\end{lemma}
\begin{widepar}
    \begin{proof}
        For the stacking operation $\mtimescatob$ on objects, we use $\mtimescatobof{\DPL}\definedas\mtimescatobof{\PosL}$.
        For stacking morphisms, we define $\mtimescatmor$ by
        % Given two \SY{design problems} % \begin{equation}
        %     \begin{aligned}
        %         \adpa\colon \posA\posop \Ptimes \posC & \toinPos \Bool, \\
        %         \adpb\colon \posB\posop \Ptimes \posD & \toinPos \Bool,
        %     \end{aligned}
        % \end{equation}
        % we define
        \begin{equation}
            \label{eq:dpmulti}
            \prfperiod{
                \adpa\colon \posAop\Ptimes \posC \mto \Bool
            }{\qquad}{
                \adpb\colon \posB\posop\Ptimes \posD \mto \Bool,
            }{
                \defmap{
                    \adpa \mtimescatmor \adpb
                }{
                    \pars{\posA \mtimescatob \posB}\op
                    \Ptimes
                    \pars{\posC \mtimescatob \posD}
                }{
                    \mto
                }{
                    \Bool
                }{
                    \tup{\maketupconcat{a^*, c^*},\ \maketupconcat{b, d}}
                }{
                    \adpa(a^*,c) \booland \adpb(c^*, d)
                }
            }
        \end{equation}
        Note that this is a valid definition of a \SY{design problem} because the expression $\adpa(a^*,c) \booland \adpb(c^*, d)$ is monotone, as may readily be checked.

        To show associativity, consider three DPs
        \begin{equation}
            \adpa\colon \posAop\Ptimes \posC \toinPos \Bool, \qquad
            \adpb\colon \posB\posop\Ptimes \posD \toinPos \Bool, \qquad
            \adpc\colon \posE\posop\Ptimes \posF \toinPos \Bool.
        \end{equation}
        We compute $(\adpa \mtimescatmor \adpb) \mtimescatmor \adpc$
        and $\adpa \mtimescatmor (\adpb \mtimescatmor \adpc)$ according to the recipe in~\cref{eq:dpmulti}:
        \begin{equation}
            \defmapcomma{
                \parslight{\adpa \mtimescatmor \adpb} \mtimescatmor \adpc
            }{
                \pars{\parslight{\posA \mtimescatob \posB} \mtimescatob \posE}\op
                \Ptimes
                \pars{\parslight{\posC \mtimescatob \posD} \mtimescatob \posF}
            }{
                \toinPos
            }{
                \Bool
            }{
                \tup{
                    \parslight{\maketupconcat{a^*, c^*}} \tupconcat e^*
                    ,\
                    \parslight{\maketupconcat{b, d}} \tupconcat f
                }
            }{
                \parslight{
                    \adpa(a^*,c) \booland \adpb(c^*, d)
                }
                \ \booland\
                \adpc(e^*, f)
            }
        \end{equation}
        \begin{equation}
            \defmapperiod{
                \adpa \mtimescatmor \parslight{\adpb \mtimescatmor \adpc}
            }{
                \pars{
                    \posA \mtimescatob
                    \parslight{
                        \posB \mtimescatob \posE
                    }
                }\op
                \Ptimes
                \pars{
                    \posC \mtimescatob
                    \parslight{
                        \posD \mtimescatob \posF
                    }
                }
            }{
                \toinPos
            }{
                \Bool
            }{
                \tup{
                    \maketupconcat{a^*, \parslight{\maketupconcat{c^*, e^*})}},\
                    \maketupconcat{b, \parslight{\maketupconcat{d, f}}}
                }
            }{
                \adpa(a^*,c)
                \ \booland\
                \parslight{
                    \adpb(c^*, d)
                    \booland
                    \adpc(e^*, f)
                }
            }
        \end{equation}
        Because the operations $\mtimescatob$ and $\booland$ are \SY{associative}, we can erase all the light parentheses in the formulas, and we find that $(\adpa \mtimescatmor \adpb) \mtimescatmor \adpc$ and
        $\adpa \mtimescatmor (\adpb \mtimescatmor \adpc)$ are equal to the design problem
        \begin{equation}
            \defmapperiod{
                \adpa \mtimescatmor \adpb \mtimescatmor \adpc
            }{
                \pars{
                    \posA \mtimescatob
                    \posB \mtimescatob \posE
                }\op
                \Ptimes
                \pars{
                    \posC \mtimescatob
                    \posD \mtimescatob \posF
                }
            }{
                \toinPos
            }{
                \Bool
            }{
                \tup{
                    \maketupconcat{a^*, c^*, e^*},\
                    \maketupconcat{b, d, f}
                }
            }{
                \adpa(a^*,c)
                \booland
                \adpb(c^*, d)
                \booland
                \adpc(e^*, f)
            }
        \end{equation}
    \end{proof}
\end{widepar}

% where
% \begin{equation}
%     (\adpa \cproddp \adpb)  (\tup{\posAnel{}\opel,\posBnel{}\opel,\posCnel{},\posDnel{}})
%     \definedas\adpa(\tup{\posAnel{}\opel,\posCnel{}})\booland
%     \adpb(\tup{\posBnel{}\opel,\posDnel{}}).
% \end{equation}

% We now consider the stacking operation on objects we considered for~\PosL, and the stacking operation on morphism defined as
% \begin{equation}
%     \defmapcomma{
%         \mtimescatmor
%     }{
%         \Mor_{\DPL} \cartprod \Mor_{\DPL}
%     }{
%         \mto
%     }{
%         \Mor_{\DPL}
%     }{
%         \tup{\adpa,\adpb}
%     }{
%         \adpa \cproddp \adpb
%     },
% \end{equation}

% \begin{lemma}\label{lem:DPL-associative-stacking}
%     \DPL, equipped with the aforementioned stacking operations on objects and morphisms, is  \SY{associative stacking}.
% \end{lemma}

% \begin{proof}
%     We show that~$\DPL$ is an  \SY{associative stacking} semicategory.
%     We first show that it is a stacking category, by first showing that the stacking operations indeed fulfill \cref{def:simple-stacking-semi-cat}, and then showing the compatibility condition.
%     The stacking operation on objects fulfills \cref{def:simple-stacking-semi-cat}, since, given two \SY{posets} of tuples it returns a \SY{poset} of tuples.
%     To show that the stacking operation on morphism fulfills \cref{def:simple-stacking-semi-cat}, we need to show that given \SY{design problems} ~$\adpa,\adpb$, the morphism~$\adpa \mtimescatmor \adpb$ is a design problem.
%     This is easy to show, by leveraging the fact that~$\adpa$ and~$\adpb$ are \SY{design problems}.
%     We know that
%     \begin{equation}
%         \prfcomma{\tup{\posAnel{}\opel,\posBnel{}\opel,\posCnel{},\posDnel{}} \posleqof{\posA\posop \Ptimes \posB\posop \Ptimes \posC\Ptimes \posD} \tup{{\posAnel{}'}\opel,{\posBnel{}'}\opel,\posCnel{}',\posDnel{}'}}
%         {(\posAnel{}\opel \posleqof{\posA\op} {\posAnel{}'}\opel) \booland (\posBnel{}\opel \posleqof{\posB\op} {\posBnel{}'}\opel)\booland (\posCnel{}\posleqof{\posC}\posCnel{}')\booland (\posDnel{}\posleqof{\posD}\posDnel{}')}
%     \end{equation}
%     and that~$\booland$ is a monotone operator.
%     Therefore,~$\adpa \cproddp \adpb$ is a design problem.
%     Furthermore, the two stacking operations are compatible:
%     \begin{equation}
%         \prfperiod{\adpa\colon \posA\posop \Ptimes \posC \toinPos \Bool}
%         {\adpb\colon \posB\posop \Ptimes \posD\toinPos \Bool}
%         {\adpa \cproddp \adpb \colon \posA\posop \Ptimes \posB\posop \Ptimes\posC\Ptimes\posD\toinPos \Bool}
%     \end{equation}

%     \begin{comment}
%     For brevity, we write~$\posAnel{}=\tup{\posAnel{1},\ldots, \posAnel{m}}$,~$\posBnel{}=\tup{\posBnel{1},\ldots, \posBnel{n}}$,~$\posCnel{}=\tup{\posCnel{1},\ldots, \posCnel{o}}$, and~$\posDnel{}=\tup{\posDnel{1},\ldots, \posDnel{t}}$.
%     We have:
%     \begin{equation}
%         \prfperiod{
%         \posAnel{}\tupconcat \posBnel{}\posleqof{\cObj{\posA_1, \ldots, \posA_m, \posB_1, \ldots, \posB_n}} \posAnel{}'\tupconcat \posBnel{}'
%         }
%         {\prftree{(\posAnel{} \posleqof{\cObj{\posA_1,\ldots, \posA_m}} \posAnel{}') \booland (\posBnel{} \posleqof{\cObj{\posB_1,\ldots, \posB_n}} \posBnel{}')}
%         {\prftree{(\adpa(\posAnel{}\opel,\posCnel{})\posgeqof{\Bool} \adpa({\posAnel{}'}\opel,\posCnel{}))
%                 \booland
%                 (\adpb(\posBnel{}\opel,\posDnel{})\posgeqof{\Bool} \adpb({\posBnel{}'}\opel,\posDnel{}))}
%             {
%                 \prftree{
%                     (\adpa({\posAnel{}}\opel,\posCnel{})\booland \adpb(\posBnel{}\opel,\posDnel{}))
%                     \posgeqof{\Bool}
%                     (\adpa({\posAnel{}'}\opel,\posCnel{})\booland \adpb({\posBnel{}'}\opel,\posDnel{}))
%                 }
%                 {(\adpa \cprod \adpb)((\posAnel{}\tupconcat \posBnel{})\opel,\posCnel{}\tupconcat \posDnel{})\posgeqof{\Bool}
%                     (\adpa \cprod \adpb)((\posAnel{}'\tupconcat \posBnel{}')\opel,\posCnel{}\tupconcat \posDnel{})} }
%             }
%         }
%     \end{equation}
%     For the second iteration, we have:
%     \begin{equation}
%         \prfperiod{
%         \posCnel{}\tupconcat \posDnel{}\posleqof{\cObj{\posC_1, \ldots, \posC_o, \posD_1, \ldots, \posD_t}} \posCnel{}'\tupconcat \posDnel{}'
%         }
%         {\prftree{(\posCnel{} \posleqof{\cObj{\posC_1,\ldots,\posC_o}} \posCnel{}') \booland (\posDnel{}\posleqof{\cObj{\posD_1,\ldots, \posD_t}} \posDnel{}')}
%         {\prftree{(\adpa(\posAnel{}\opel,\posCnel{})\posleqof{\Bool}
%                 \adpa(\posAnel{}\opel,\posCnel{}'))
%                 \booland
%                 (\adpb(\posBnel{}\opel,\posDnel{})\posleqof{\Bool}
%                 \adpb(\posBnel{}\opel,\posDnel{}'))}
%             {
%                 \prftree{
%                     (\adpa(\posAnel{}\opel,\posCnel{})\booland \adpb(\posBnel{}\opel,\posDnel{}))
%                     \posleqof{\Bool}
%                     (\adpa(\posAnel{}\opel,\posCnel{}')\booland \adpb(\posBnel{}\opel,\posDnel{}'))
%                 }
%                 {(\adpa \cprod \adpb)((\posAnel{}\tupconcat \posBnel{})\opel,\posCnel{}\tupconcat \posDnel{})
%                     \posleqof{\Bool}
%                     (\adpa \cprod \adpb)((\posAnel{}\tupconcat \posBnel{})\opel,\posCnel{}'\tupconcat \posDnel{}')} }
%             }
%         }
%     \end{equation}
%     Therefore,~$\adpa \cprod \adpb$ is really a design problem.
%     \end{comment}
%     Associativity can be shown following what we did for~$\cObj{\Set}$.
% \end{proof}





\begin{example}[\SetL is a functorial stacking semicategory]
    \label{ex:setfunstack}
    We want to show that~\SetL is a \SY{functorial stacking semicategory}.

    Consider four morphisms
    \begin{equation}
        \begin{aligned}
            \mora\colon \setA \mto \setB, & \quad \morc\colon \setB \mto \setC, \\
            \morb\colon \setD \mto \setE, & \quad \mord\colon \setE \mto \setF.
            \\
        \end{aligned}
    \end{equation}
    We want to show that
    \begin{equation}
        (\mora \mtimescatmor \morb )
        \mthen (\morc \mtimescatmor \mord )
        =
        (\mora \mthen \morc ) \mtimescatmor (\morb \mthen \mord ),
    \end{equation}
    We show this by showing that, for any $\setAel\tupconcat \setDel\setin \setA\cprod\setD$:
    \begin{equation}
        \begin{aligned}
            ((\mora \mtimescatmor \morb )\mthen (\morc \mtimescatmor \mord ))(\setAeln{}\tupconcat \setDeln{})
             & = (\morc \mtimescatmor \mord)(\mora(\setAeln{})\tupconcat \morb(\setDeln{})) \\
             & =\morc(\mora(\setAeln{}))\tupconcat \mord(\morb(\setDeln{})) \\
             & =(\mora \mthen \morc)(\setAeln{})\tupconcat (\morb\mthen \mord)(\setDeln{}) \\
             & =((\mora \mthen \morc ) \mtimescatmor (\morb \mthen \mord ))(\setAeln{}\tupconcat \setDeln{}).
        \end{aligned}
    \end{equation}
\end{example}

\begin{exercise}
    \label{ex:the-usual-suspects-are-functorial-stacking-categories}
    Prove that \PosL is a functorial stacking category.
    Hint: You have to define the functor $\mtimescat$ and check that it satisfies the two equations in \cref{def:functorial-stacking-cat}.
\end{exercise}
\begin{solution}
    \todotext{Solution for \cref{ex:the-usual-suspects-are-functorial-stacking-categories}}
\end{solution}


\begin{gradedexercise}[\exname{RelFunStack}]
    \label{ex:RelFunStack}
    Prove that the structure defined in \cref{ex:relassocstack} makes~\RelL a \SY{functorial stacking semicategory}.
\end{gradedexercise}

\solutionof{RelFunStack}

\begin{lemma}[\DPL is a functorial stacking semicategory]\label{lem:DPL-functorial-stacking}
    \DPL, equipped with the aforementioned stacking operations on objects and morphisms, is \SY{functorial stacking semicategory}.
\end{lemma}

\begin{proof}
    Consider
    \begin{equation}
        \begin{aligned}
            \adpa\colon \posAop\Ptimes \posC     & \toinPos \Bool \\
            \adpb\colon \posB\posop\Ptimes \posD & \toinPos \Bool \\
            \adpc\colon \posC\posop\Ptimes \posE & \toinPos \Bool \\
            \adpd\colon \posD\posop\Ptimes \posF & \toinPos \Bool
        \end{aligned}
    \end{equation}
    We want to prove that
    \begin{equation}
        (\adpa \dpthen \adpc)
        \mtimescatmor (\adpb \dpthen \adpd)=(\adpa \mtimescatmor \adpb) \dpthen (\adpc \mtimescatmor \adpd).
    \end{equation}
    %For brevity, in the following we write
    %\begin{align*}
    %    \posAnel{} & =\tup{\posAnel{1},\ldots, \posAnel{i}}, &
    %    \posBnel{} & =\tup{\posBnel{1},\ldots, \posBnel{l}} \\
    %    \posCnel{} & =\tup{\posCnel{1},\ldots, \posCnel{m}}, &
    %    \posDnel{} & =\tup{\posDnel{1},\ldots, \posDnel{n}} \\
    %    \posEnel{} & =\tup{\posEnel{1},\ldots, \posEnel{o}}, &
    %    \posFnel{} & =\tup{\posFnel{1},\ldots, \posFnel{v}}
    %\end{align*}
    We start from the left-hand side.
    We have
    \begin{equation}
        (\adpa \dpthen \adpc)(\posAnel{}\opel,\posEnel{})
        =\bigvee_{\posCnel{}\setin \posC}
        \adpa(\posAnel{}\opel, \posCnel{}) \booland \adpc(\posCnel{}\opel,\posEnel{})
    \end{equation}
    and
    \begin{equation}
        \begin{aligned}
            (\adpb \dpthen \adpd)(\posBnel{}\opel,\posFnel{})
            =\bigvee_{\posDnel{}\setin \posD}
            \adpb(\posBnel{}\opel, \posDnel{}) \booland \adpd(\posDnel{}\opel\posFnel{})
        \end{aligned}
    \end{equation}
    Therefore, we know
    \begin{equation}
        \begin{aligned}
             & ((\adpa \dpthen \adpc)\mtimescatmor (\adpb \dpthen \adpd))
            ((\posAnel{}\tupconcat \posBnel{})\opel,\posEnel{}\tupconcat \posFnel{}) \\
             & =\bigvee_{\posCnel{}\setin \posC}
            \adpa(\posAnel{}\opel, \posCnel{}) \booland \adpc(\posCnel{}\opel,\posEnel{}) \booland
            \bigvee_{\posDnel{}\setin \posD} \adpb(\posBnel{}\opel, \posDnel{}) \booland \adpd(\posDnel{}\opel,\posFnel{}).
        \end{aligned}
    \end{equation}

    On the other hand, we have
    \begin{equation}
        (\adpa \mtimescatmor \adpb)
        ((\posAnel{}\tupconcat \posBnel{})\opel, \posCnel{}\tupconcat \posDnel{})
        =\adpa(\posAnel{}\opel,\posCnel{}) \booland \adpb(\posBnel{}\opel,\posDnel{})
    \end{equation}
    and
    \begin{equation}
        (\adpc \mtimescatmor \adpd)
        ((\posCnel{}\tupconcat \posDnel{})\opel, \posEnel{}\tupconcat \posFnel{})
        =\adpc(\posCnel{}\opel,\posEnel{}) \booland \adpd(\posDnel{}\opel,\posFnel{})
    \end{equation}
    Therefore, we know
    \begin{equation}
        \begin{aligned}
             & ((\adpa \mtimescatmor \adpb)\dpthen (\adpc \mtimescatmor \adpd))((\posAnel{}\tupconcat \posBnel{})\opel,\posEnel{}\tupconcat \posFnel{}) \\
             & =\bigvee_{\posCnel{}\tupconcat \posDnel{}\setin \cObj{\posC,\posD}}
            (\adpa \mtimescatmor \adpb)((\posAnel{}\tupconcat \posBnel{})\opel, \posCnel{}\tupconcat \posDnel{})
            \booland (\adpc \mtimescatmor \adpd)((\posCnel{}\tupconcat \posDnel{})\opel, \posEnel{}\tupconcat \posFnel{}) \\
             & =\bigvee_{\posCnel{}\tupconcat \posDnel{}\setin \cObj{\posC,\posD}}
            \adpa(\posAnel{}\opel,\posCnel{}) \booland \adpb(\posBnel{}\opel,\posDnel{})
            \booland
            \adpc(\posCnel{}\opel,\posEnel{}) \booland \adpd(\posDnel{}\opel,\posFnel{}) \\
             & =\bigvee_{\posCnel{}\setin \posC}
            \adpa(\posAnel{}\opel, \posCnel{}) \booland \adpc(\posCnel{}\opel,\posEnel{}) \booland
            \bigvee_{\posDnel{}\setin \posD} \adpb(\posBnel{}\opel, \posDnel{}) \booland \adpd(\posDnel{}\opel,\posFnel{}),
        \end{aligned}
    \end{equation}
    proving the statement for any \SY{posets}~$\posA,\posB,\posC,\posD,\posE,\posF$ (and hence, also for \SY{posets} of tuples).

\end{proof}
